% ============================================================
\subsubsection{Desarrollo de la interfaz del sistema}
% ============================================================

\paragraph{Tecnologías}

Las Tablas~\ref{tab:tec_back} y~\ref{tab:tec_front} presentan las tecnologías
empleadas y sus versiones, tomadas de los archivos \texttt{package.json} del
proyecto.

\begin{table}[H]
\centering\small
\caption{Tecnologías y librerías principales del backend}\label{tab:tec_back}
\begin{tabular}{|l|l|l|}
\hline
\textbf{Librería} & \textbf{Versión} & \textbf{Propósito} \\ \hline
Node.js & \pendiente{v} & Entorno de ejecución \\ \hline
TypeScript & \^{}5.4.5 & Tipado estático \\ \hline
Express & \^{}4.22.2 & API REST \\ \hline
Prisma (\texttt{@prisma/client}) & \^{}5.14.0 & ORM sobre PostgreSQL \\ \hline
PostgreSQL & \pendiente{v} & Base de datos relacional \\ \hline
openai (SDK, apuntando a Groq) & \^{}6.37.0 & Cliente de la API de inferencia \\ \hline
pdf-parse & \^{}2.4.5 & Extracción de texto de PDF \\ \hline
multer & \^{}2.2.0 & Recepción de archivos \\ \hline
jsonwebtoken & \^{}9.0.2 & Tokens de sesión JWT \\ \hline
bcryptjs & \^{}2.4.3 & Cifrado de contraseñas \\ \hline
helmet & \^{}8.2.0 & Cabeceras de seguridad HTTP \\ \hline
\end{tabular}
\end{table}

\begin{table}[H]
\centering\small
\caption{Tecnologías y librerías principales del frontend}\label{tab:tec_front}
\begin{tabular}{|l|l|l|}
\hline
\textbf{Librería} & \textbf{Versión} & \textbf{Propósito} \\ \hline
React & \^{}18.3.1 & Construcción de la interfaz \\ \hline
Vite & \^{}5.3.1 & Empaquetado y servidor de desarrollo \\ \hline
TailwindCSS & \^{}3.4.4 & Estilos \\ \hline
axios & \^{}1.7.2 & Consumo de la API REST \\ \hline
react-markdown + remark-gfm & \^{}10.1.0 / \^{}4.0.1 & Renderizado de respuestas del asistente con tablas \\ \hline
lucide-react & \^{}1.48.0 & Iconografía \\ \hline
\end{tabular}
\end{table}

\paragraph{Estructura del proyecto}

El backend sigue una organización por capas: \texttt{routes} define los
endpoints, \texttt{controllers} contiene la lógica de cada recurso,
\texttt{services} encapsula la integración con el modelo de lenguaje,
\texttt{middleware} gestiona autenticación y errores, y \texttt{prisma} contiene
el esquema, las migraciones y la carga inicial del catálogo
(Figura~\ref{fig:estructura_back}). El frontend organiza sus componentes por
módulo funcional en \texttt{components/features} (pacientes, historia clínica,
antropometría, biometría, hábitos dietéticos, alimentos, menús, progreso y
administración) (Figura~\ref{fig:estructura_front}).

\figura[0.4\linewidth]{estructura_backend.png}{Estructura del backend}{fig:estructura_back}
\figura[0.4\linewidth]{estructura_frontend.png}{Estructura del frontend}{fig:estructura_front}

\paragraph{Modelo de base de datos}

El esquema de Prisma define diez entidades: \texttt{User}, \texttt{Patient},
\texttt{ClinicalHistory}, \texttt{Biometrics}, \texttt{Anthropometry},
\texttt{DietaryHabits}, \texttt{ChatMessage}, \texttt{Food},
\texttt{WeeklyMenu} y \texttt{WeeklyMenuItem} (Figura~\ref{fig:er}).

\figura{diagrama_er.png}{Modelo entidad-relación de la base de datos}{fig:er}

\paragraph{Trazabilidad de requisitos en la implementación}

La Tabla~\ref{tab:traza_impl} relaciona cada requisito funcional con el
componente que lo implementa.

\begin{table}[H]
\centering\footnotesize
\caption{Requisitos funcionales y componente que los implementa}\label{tab:traza_impl}
\begin{tabularx}{\linewidth}{|l|X|X|}
\hline
\textbf{RF} & \textbf{Endpoint} & \textbf{Componente de interfaz} \\ \hline
RF01 & \texttt{POST/PUT /api/patients} & \texttt{features/patients} \\ \hline
RF02 & \texttt{PUT /api/patients/:id/clinical-history} & \texttt{features/clinical} \\ \hline
RF03 & \texttt{POST /api/patients/:id/anthropometry} & \texttt{features/anthropometry} \\ \hline
RF04 & --- (cálculo en cliente) & \texttt{utils/nutritionCalculations.ts} \\ \hline
RF05 & \texttt{POST /api/patients/:id/biometrics} & \texttt{features/biometrics} \\ \hline
RF06 & \texttt{POST /api/patients/:id/dietary-habits} & \texttt{features/dietary} \\ \hline
RF07 & --- (cálculo en cliente) & \texttt{NutritionistPanel.tsx} \\ \hline
RF08 & --- (cálculo en cliente) & \pendiente{componente del gráfico de dona} \\ \hline
RF09 & \texttt{POST /api/ai/chat} (con \texttt{patientId}) & \texttt{NutritionistPanel.tsx} \\ \hline
RF10 & \texttt{POST /api/ai/chat} (global) & \texttt{NutritionistPanel.tsx} \\ \hline
RF11 & \texttt{POST /api/ai/chat-upload} & \texttt{NutritionistPanel.tsx} \\ \hline
RF12 & \texttt{GET/POST/PUT/DELETE /api/foods} & \texttt{features/foods} \\ \hline
\end{tabularx}
\end{table}

\paragraph{Interfaz del sistema}

Las Figuras~\ref{fig:ui_login} a~\ref{fig:ui_pdf} muestran los módulos
principales del prototipo en funcionamiento.

\figura{im02_login.png}{Pantalla de inicio de sesión}{fig:ui_login}
\figura{im01_dashboard.png}{Panel principal del nutricionista}{fig:ui_dashboard}
\figura{im03_pacientes.png}{Listado de pacientes}{fig:ui_pacientes}
\figura{im04_registro_paciente.png}{Registro de paciente}{fig:ui_registro}
\figura{im05_antropometria.png}{Registro de medidas antropométricas con cálculo de IMC e ICC}{fig:ui_antropo}
\figura{im06_biometria.png}{Registro de datos biométricos}{fig:ui_bio}
\figura{im07_chat_ia.png}{Asistente de IA con contexto del paciente}{fig:ui_chat}
\figura{im08_chat_pdf.png}{Consulta al asistente con un documento PDF}{fig:ui_pdf}

\pendiente{añade capturas del cálculo de GET, del gráfico de macronutrientes
(RF07, RF08) y del catálogo de alimentos (RF12) si existen en la interfaz}

El código fuente se encuentra en el repositorio
\pendiente{URL del repositorio de GitHub}.
